% !TeX root = ../../elegantbook-cn.tex

\begin{example}
（2026东城期末）设函数 $f(x)=\mathrm{e}^{ax}-\mathrm{e}^{bx}$．\\
（1）当 $a=2,b=1$ 时，求曲线 $y=f(x)$ 在点 $(0,f(0))$ 处的切线方程；\\
（2）当 $a=3,b=2$ 时，\\
（i）求 $f(x)$ 的极值；\\
（ii）若 $\lambda\in\left[2,\dfrac52\right]$，求证：关于 $x$ 的方程 $\dfrac{f(x)}{x}=\mathrm{e}^{\lambda x}$ 在 $(0,+\infty)$ 上无解．
\end{example}

\begin{example}
（2025东城期末）设函数 $f(x)=\dfrac{\ln(x+1)}{\mathrm{e}^x}$．\\
（1）求曲线 $y=f(x)$ 在点 $(0,f(0))$ 处的切线方程；\\
（2）求 $f(x)$ 的极值点个数；\\
（3）若 $x>0$ 时，$\dfrac{f(x)}{x}>k$，求 $k$ 的取值范围．
\end{example}

\begin{example}
（2024东城期末）已知函数 $f(x)=\dfrac{x-1}{x+1}-k\mathrm{e}^x$，$k>0$．\\
（1）若 $k=1$，求曲线 $y=f(x)$ 在点 $(0,f(0))$ 处的切线方程；\\
（2）若 $1\leqslant k<2$，证明：函数 $y=f(x)$ 在 $(0,+\infty)$ 上有极大值 $m$，且 $-3<m<1$．
\end{example}

\begin{example}
（2023东城期末）已知函数 $f(x)=x\mathrm{e}^x$．\\
（1）求曲线 $y=f(x)$ 在点 $(0,f(0))$ 处的切线方程；\\
（2）求 $f(x)$ 的极值；\\
（3）证明：当 $m\leqslant1$ 时，曲线 $C_1:y=f(x)$ 与曲线 $C_2:y=\ln x+x+m$ 至多存在一个交点．
\end{example}

\begin{example}
（2022东城期末）曲线 $y=\ln x$ 在点 $A(t,\ln t)$ 处的切线 $l$ 交 $x$ 轴于点 $M$．\\
（1）当 $t=\mathrm{e}$ 时，求切线 $l$ 的方程；\\
（2）$O$ 为坐标原点，记 $\triangle AMO$ 的面积为 $S$，求面积 $S$ 以 $t$ 为自变量的函数解析式，写出其定义域，并求单调区间．
\end{example}

\begin{example}
（2021东城期末）已知函数 $f(x)=1-\dfrac{ax^2}{\mathrm{e}^x}$，$a\in\mathbb{R}$．\\
（1）若曲线 $y=f(x)$ 在点 $(1,f(1))$ 处的切线平行于直线 $y=x$，求该切线方程；\\
（2）若 $a=1$，求证：当 $x>0$ 时，$f(x)>0$；\\
（3）若 $f(x)$ 有且只有两个零点，求 $a$ 的值．
\end{example}

\begin{example}
（2026西城期末）已知函数 $f(x)=(2x-a)\mathrm{e}^x$ 满足 $f'(1)-3f(1)=0$，其中 $a\in\mathbb{R}$．\\
（1）求 $a$ 的值；\\
（2）证明：在曲线 $y=f(x)$ 的所有切线中，有且仅有一条切线与直线 $y=2x$ 平行；\\
（3）若存在 $m$，使得对于任意 $x_1,x_2\in(0,+\infty)$，且 $x_1<x_2$，都有 $|f(x_2)-f(x_1)|>m(x_2^2-x_1^2)$，求实数 $m$ 的取值范围．
\end{example}

\begin{example}
（2025西城期末）已知函数 $f(x)=\ln(ax+1)-x$，其中 $a>0$．\\
（1）当 $a=1$ 时，求曲线 $y=f(x)$ 在点 $(0,f(0))$ 处的切线方程；\\
（2）当 $a=2$ 时，证明：对任意 $t\in(0,+\infty)$，曲线 $y=f(x)$ 总在直线 $y=x+t$ 的下方；\\
（3）若函数 $f(x)$ 有两个零点 $x_1,x_2$，且 $0<x_2-x_1<1$，求 $a$ 的取值范围．
\end{example}

\begin{example}
（2024西城期末）已知函数 $f(x)=\dfrac{\mathrm{e}^{ax}}{x}$，其中 $a>0$．\\
（1）当 $a=1$ 时，求曲线 $y=f(x)$ 在点 $(1,f(1))$ 处的切线方程；\\
（2）求 $f(x)$ 的单调区间；\\
（3）当 $x_1<x_2$ 且 $x_1x_2>0$ 时，判断 $f(x_1)-f(x_2)$ 与 $\dfrac1{x_1}-\dfrac1{x_2}$ 的大小，并说明理由．
\end{example}

\begin{example}
（2023西城期末）已知函数 $f(x)=a\ln x+x\mathrm{e}^x-\mathrm{e}$，其中 $a\in\mathbb{R}$．\\
（1）当 $a=0$ 时，求曲线 $y=f(x)$ 在点 $(1,f(1))$ 处的切线方程；\\
（2）当 $a>0$ 时，判断 $f(x)$ 的零点个数，并加以证明；\\
（3）当 $a<0$ 时，证明存在实数 $m$，使 $f(x)\geqslant m$ 恒成立．
\end{example}
\begin{example}
（2022西城期末）已知函数 $f(x)=\mathrm{e}^x(x^2+ax+1)$．\\
（I）若 $a=0$，求 $f(x)$ 在点 $(0,f(0))$ 处的切线方程；\\
（II）若 $f(x)$ 在 $(-1,1)$ 上恰有一个极小值点，求实数 $a$ 的取值范围；\\
（III）若对于任意 $x\in\left(0,\dfrac{\pi}{2}\right]$，$f(x)>\mathrm{e}^x(x^2\cos x+1)$ 恒成立，求实数 $a$ 的取值范围．
\end{example}


\begin{example}
（2021西城期末）已知函数 $f(x)=x^3-x$．\\
（I）求曲线 $y=f(x)$ 在点 $(1,f(1))$ 处的切线方程；\\
（II）求函数 $f(x)$ 的单调区间和极值；\\
（III）设函数 $t(x)=\dfrac{f(x)}{x\sin x}-2$，$x\in(0,\pi)$，试判断 $t(x)$ 的零点个数，并证明你的结论．
\end{example}

\begin{example}
（2026海淀期末）已知函数
\[
f(x)=
\begin{cases}
x^3-3ax+a,&x<0,\\
(x^2+a)\mathrm{e}^x,&x\geqslant0.
\end{cases}
\]
（1）当 $a=1$ 时，求曲线 $y=f(x)$ 在点 $(-2,f(-2))$ 处的切线在 $y$ 轴上的截距；\\
（2）若函数 $f(x)$ 在 $(-\infty,+\infty)$ 上单调递增，求 $a$ 的值；\\
（3）若函数 $f(x)$ 在 $x=s$ 处取得极小值，求 $\dfrac{s}{f(s)}$ 的取值范围．
\end{example}

\begin{example}
（2025海淀期末）已知函数 $f(x)=\dfrac{\ln(ax)}{x-1}$．\\
（1）当 $a=1$ 时，求 $f(x)$ 的定义域；\\
（2）若 $f(x)$ 在区间 $(-\infty,-1]$ 上单调递减，求 $a$ 的取值范围；\\
（3）当 $a=\dfrac{2}{\mathrm{e}}$ 时，证明：若 $x_1\in(0,1)$，$x_2\in(1,+\infty)$，则 $f(x_1)-f(x_2)>\dfrac32$．（参考数据：$\mathrm{e}^2\approx7.39$，$\mathrm{e}^3\approx20.09$，$\mathrm{e}^4\approx54.60$）
\end{example}

\begin{example}
（2024海淀期末）已知函数 $f(x)=ax^2-x\sin x+b$．\\
（1）当 $a=1$ 时，求证：\\
①当 $x>0$ 时，$f(x)>b$；\\
②函数 $f(x)$ 有唯一极值点；\\
（2）若曲线 $C_1$ 与曲线 $C_2$ 在某公共点处的切线重合，则称该切线为 $C_1$ 和 $C_2$ 的“优切线”．若曲线 $y=f(x)$ 与曲线 $y=-\cos x$ 存在两条互相垂直的“优切线”，求 $a,b$ 的值．
\end{example}

\begin{example}
（2023海淀期末）已知函数 $f(x)=x\ln(x+1)$．\\
（1）判断 $0$ 是否为 $f(x)$ 的极小值点，并说明理由；\\
（2）证明：$\dfrac{f(x)}{x^2}>-\dfrac12x+1$．
\end{example}

\begin{example}
（2022海淀期末）函数 $f(x)=a\mathrm{e}^x-\sin x+2x$．\\
（1）求曲线 $y=f(x)$ 在点 $(0,f(0))$ 处的切线方程；\\
（2）当 $a\geqslant0$ 时，求函数 $f(x)$ 在 $[0,1]$ 上的最小值；\\
（3）直接写出 $a$ 的一个值，使 $f(x)\leqslant a$ 恒成立，并证明．
\end{example}

\begin{example}
（2021海淀期末）已知函数 $f(x)=\dfrac{\ln x}{x}$．\\
（1）求函数 $f(x)$ 的单调区间；\\
（2）设 $g(x)=f(x)-x$，求证：$g(x)\leqslant-1$；\\
（3）设 $h(x)=f(x)-x^2+2ax-4a^2+1$，若存在 $x_0$ 使得 $h(x_0)\geqslant0$，求 $a$ 的最大值．
\end{example}

\begin{example}
（2026丰台期末）已知函数 $f(x)=\ln(ax+1)-ax^2$（$a\neq0$）．\\
（1）若 $x=\dfrac14$ 为 $f(x)$ 的极值点，求实数 $a$ 的值；\\
（2）若 $a<0$，讨论 $f(x)$ 的单调区间；\\
（3）若 $a>0$，证明：当 $x\neq0$ 时，曲线 $y=f(x)$ 在 $(0,f(0))$ 处的切线总在曲线 $y=f(x)$ 的上方．
\end{example}

\begin{example}
（2025丰台期末）已知函数 $f(x)=\mathrm{e}^x-x$．\\
（1）求曲线 $y=f(x)$ 在点 $(0,f(0))$ 处的切线方程；\\
（2）求 $f(x)$ 的极值；\\
（3）设函数 $g(x)=f(x)+x^2+x$，求证：$g(x)$ 的最小值大于 $\dfrac12$．
\end{example}

\begin{example}
（2024丰台期末）已知函数 $f(x)=\mathrm{e}^x(x^2-ax-a)$．\\
（1）若曲线 $y=f(x)$ 在点 $(1,f(1))$ 处的切线平行于 $x$ 轴，求实数 $a$ 的值；\\
（2）求函数 $f(x)$ 的单调区间．
\end{example}

\begin{example}
（2023丰台期末）已知函数 $f(x)=\ln x+\sin x$．\\
（1）求曲线 $y=f(x)$ 在点 $(1,f(1))$ 处的切线方程；\\
（2）求函数 $f(x)$ 在区间 $[1,\mathrm{e}]$ 上的最小值；\\
（3）证明函数 $f(x)$ 只有一个零点．
\end{example}

\begin{example}
（2022丰台期末）已知函数 $f(x)=x^2-a\ln x$（$a\in\mathbb{R}$ 且 $a\neq0$）．\\
（1）当 $a=1$ 时，求曲线 $y=f(x)$ 在点 $(1,f(1))$ 处的切线方程；\\
（2）若 $f(x)\geqslant0$ 恒成立，求 $a$ 的取值范围．
\end{example}

\begin{example}
（2021丰台期末）已知函数 $f(x)=(x-a)\mathrm{e}^x$（$a\in\mathbb{R}$）．\\
（1）当 $a=1$ 时，求曲线 $y=f(x)$ 在点 $(1,f(1))$ 处的切线方程；\\
（2）如果函数 $f(x)$ 在区间 $(0,1)$ 上有极值，且 $f(x)+a\leqslant0$ 对于 $x\in[0,1]$ 恒成立，求 $a$ 的取值范围．
\end{example}

\begin{example}
（2026朝阳期末）已知函数 $f(x)=x\ln(x+a)$（$a>0$），曲线 $y=f(x)$ 在点 $(0,f(0))$ 处的切线方程为 $y=0$．\\
（1）求 $a$ 的值；\\
（2）求 $f(x)$ 的单调区间；\\
（3）若 $0<m<1$，曲线 $y=f(x)$ 在点 $A(m,f(m))$、$B(-m,f(-m))$ 处的切线分别与 $x$ 轴交于 $C(x_1,0)$、$D(x_2,0)$，求证：$x_1+x_2<0$．
\end{example}

\begin{example}
（2025朝阳期末）已知函数 $f(x)=\dfrac{x^n}{\mathrm{e}^{x-1}}$（$x>0$），其中 $n$ 是常数，$n\in\mathbb{Z}$．\\
（1）当 $n=0$ 时，求曲线 $y=f(x)$ 在点 $(1,f(1))$ 处的切线方程；\\
（2）求 $f(x)$ 的极值．
\end{example}

\begin{example}
（2024朝阳期末）已知函数 $f(x)=x-a\ln x-1$（$a\in\mathbb{R}$）．\\
（1）若曲线 $y=f(x)$ 在点 $(1,0)$ 处的切线为 $x$ 轴，求 $a$ 的值；\\
（2）讨论 $f(x)$ 在区间 $(1,+\infty)$ 内极值点的个数；\\
（3）若 $f(x)$ 在区间 $(1,+\infty)$ 内有零点 $t$，求证：$t<a^2$．
\end{example}

\begin{example}
（2023朝阳期末）已知函数 $f(x)=\dfrac{\ln x}{ax}$（$a>0$）．\\
（1）求 $f(x)$ 的单调区间；\\
（2）若 $f(x)\leqslant x-\dfrac1a$ 对 $x\in(0,+\infty)$ 恒成立，求 $a$ 的取值范围；\\
（3）若 $x_2\ln x_1+x_1\ln x_2=0$（$x_1\neq x_2$），证明：$x_1+x_2>2$．
\end{example}

\begin{example}
（2022朝阳期末）已知函数 $f(x)=2\ln x-x-\ln a$，$a>0$．\\
（1）求曲线 $y=f(x)$ 在 $(1,f(1))$ 处切线的斜率；\\
（2）求函数 $f(x)$ 的极大值；\\
（3）设 $g(x)=a\mathrm{e}^x-x^2$，当 $a\in(1,\mathrm{e})$ 时，求函数 $g(x)$ 的零点个数，并说明理由．
\end{example}

\begin{example}
（2021朝阳期末）已知函数 $f(x)=\ln x-(a+2)x+ax^2$（$a\in\mathbb{R}$）．\\
（1）当 $a=0$ 时，求曲线 $y=f(x)$ 在点 $(1,f(1))$ 处的切线方程；\\
（2）求 $f(x)$ 的单调区间；\\
（3）若 $f(x)$ 恰有两个零点，求实数 $a$ 的取值范围．
\end{example}
